\documentclass{article}
\usepackage[utf8]{inputenc}
\usepackage[T1]{fontenc}
\usepackage[french]{babel}

\title{La bible de la marijuana}
\author{Bob Marley}
\date{6/02/1945}

\begin{document}

\maketitle
\newpage
\tableofcontents
\newpage

\section{Culture et séchage}
\begin{enumerate}
    \item L'art de la culture hydroponique
    \begin{enumerate}
        \item Je suis la sous-liste
        \item \label{ref_sous_liste} Le deuxième élément de la sous-liste
        \item sur le poisson et le pêcheur
        \item Le poisson mange le pêcheur
        \item Tu t'y attendais pas
    \end{enumerate}                                                          
    \item \label{ref_num2} Le séchage digne d'un sèche-cheveux (référence à la sous-liste : \ref{ref_sous_liste})
    \item Une cuisson à vous faire partie du crew d'Apollo 13
\end{enumerate}

\section{Musique et ambiance}
\begin{itemize}
    \item L'art du beat sur les nuages (Référence au deuxième élément de la liste numérotée : \ref{ref_num2})
    \item Quand tu clappes clappes clappes c'est ta façon d'aimer (Référence au deuxième élément de la liste numérotée : \ref{ref_num2})
    \item Bouge ta tête et pas ton boule (Référence au deuxième élément de la liste numérotée : \ref{ref_num2})
\end{itemize}

\section{Description et détails}
\begin{description}
	\item[Hydroponique] Référence au deuxième élément de la sous-liste numérotée : \ref{ref_sous_liste}.
    \item[Séchage] Référence au deuxième élément de la sous-liste numérotée : \ref{ref_sous_liste}.
    \item[Cuisson] Référence au deuxième élément de la sous-liste numérotée : \ref{ref_sous_liste}.
\end{description}

\end{document}